% Gerado por tools/tei2latex.py a partir de vol2-theodicea.existencia-de-deus.xml. Regenerar, não editar.

%%>>> parte-2

%%>>> cap-6

\chapter*{CAPITULO PRIMEIRO\\[0.3em]\large EXISTENCIA DE DEUS}
\addcontentsline{toc}{chapter}{CAPITULO PRIMEIRO --- EXISTENCIA DE DEUS}

{\small\noindent\textsc{Summario}: --- Demonstrabilidade da existencia de \textsc{Deus}. --- Demonstração da existencia de \textsc{Deus}. --- A existencia de \textsc{Deus} e o atheismo.\par}

%%>>> art-36

\section*{Artigo I.\\ Demonstrabilidade da existencia de Deus.}
\addcontentsline{toc}{section}{Artigo I. --- Demonstrabilidade da existencia de Deus.}

\pg{421} 6. A \textbf{existencia de Deus} póde ser \textbf{demonstrada pela razão humana}. --- A razão humana póde deduzir da existencia dos effeitos a existencia da causa. Ora todos os seres sensiveis são effeitos. Logo a razão póde deduzir da existencia dos seres sensiveis a existencia da sua causa, que é \textsc{Deus} \footnote{Os systemas, que contestam a nossa proposição são --- o \emph{tradicionalismo}, admittindo que a existencia de \textsc{Deus} só pode ser conhecida pela\pgn{422} \emph{revelação} ou pela \emph{tradição}, --- o \emph{Kantismo}, sustentando que percebemos os \emph{phenomenos} ou as \emph{manifestações}, mas não os \emph{noumenos} ou as \emph{substancias}, --- o \emph{positivismo} e o \emph{sensualismo}, affirmando que a razão humana só attinge objectos sensiveis, --- o \emph{agnosticismo}, ensinando que toda a nossa sciencia se resume e concentra na esphera dos phenomenos e que por isso o \emph{ultraphenomenico} é \emph{incognoscivel}, --- o \emph{immanentismo}, negando á razão a possibilidade de demonstrar a existencia de Deus e attribuindo-a á \emph{consciencia}.\par Por se sufficientemente conhecida a natureza dos primeiros systemas apontados, falta-nos agora explicar a natureza do ultimo, que é o \emph{immanentismo}. Segundo as pretenções d'este systema, a razão não pode demonstrar a existencia de \textsc{Deus}, como não pode demonstrar nenhuma verdade religiosa e moral, porque esta demonstração não pode ser verificada pelos dados da experiencia. As verdades religiosas e moraes não se provam com os raciocinios, como se provam as verdades mathematicas; devem brotar do fundo da nossa alma, do intimo da nossa consciencia. Se, pois, quizermos encontral-as, devemos reflectir sobre os nossos actos, investigar as exigencias do nosso espirito, sondar as necessidades do nosso coração, escutar a voz da consciencia, onde residem os dictames religiosos e moraes. Para esse trabalho é indispensavel uma vontade bóa e recta. Por quanto, os sentimentos religiosos e moraes, embora se encontrem em todos os homens, mesmo nos incredulos, comtudo só nos homens de bóa vontade e de vida irreprehensivel se encontram num estado \emph{explicito}, \emph{claro} e \emph{consciente}. --- Esta é a chamada \emph{philosophia nova}, ou philosophia da \emph{immanencia}, da \emph{acção}, do \emph{voluntarismo}, do \emph{pragmatismo}, que tem por fundadores \emph{Renouvier}, \emph{Boutroux}, \emph{Bergson}, \emph{Mach}, \emph{William James}. E' uma emanação da doutrina de \emph{Kant}, que, depois de ter negado todo o valor á razão \emph{especulativa}, quando esta investiga o \emph{sobresensivel}, ensina que a razão \emph{pratica} deve admittir a existencia de mesmo \emph{sobresensivel}, como postulado ou presupposto necessario para a ordem moral. --- Esta theoria tem sido perfilhada pelos \emph{modernistas}. Dizem elles que a existencia de Deus não pode ser demonstrada por meio dos factos \emph{externos}, isto é, pela existencia e natureza dos seres \emph{creados}, mas só por meio dos factos \emph{internos} e \emph{subjectivos}. A razão não pode demonstrar a existencia de Deus, ou não o pode fazer sem a cooperação da vontade, do sentimento e dos outros factores \emph{subjectivos}, que \emph{supprem a deficencia dos argumentos racionaes}.\par Tendo refutado na \emph{Anthropologia} o \emph{tradicionalismo}, o \emph{Kantismo}, o \emph{positivismo} e o \emph{sensualismo}, resta-nos refutar o \emph{agnosticismo} e o \emph{immanentismo}.\par A pretenção do \emph{agnosticismo}, que ensina ser \emph{incognoscivel} para a razão a existencia de Deus, é uma pretenção absurda. Que as sciencias \emph{experimentaes} não reconheçam senão os \emph{phenomenos}, e que para ellas o \emph{sobresensivel} seja \emph{incognoscivel}, concedemol-o de bom grado. Mas seguir-se-ha d'ahi que, alem do sensivel, não existe mais nada, e que, atraz do phenomeno, não existe o \emph{noumeno}? Não, com certeza. Se o sobresensivel não é objecto das sciencias experimentaes, é objecto de outra sciencia mais elevada. Os adversarios imitam nisso o celebre D. Ferrante, nos \emph{promessos esposos} de Manzoni, que negava a existencia da peste em Milão, só porque esta não cabia em nenhuma das suas categorias! Se a existencia de Deus não é perceptivel pelos\pgn{423} sentidos, é perceptivel pela razão. E a razão não prescinde dos dados da experiencia sensivel, mas funda-se nelles, explica-os e completa-os. Por quanto, no fundo de cada facto sensivel está sempre uma questão \emph{metaphysica}, que diz respeito á \emph{origem} e á \emph{natureza} do proprio facto e que só a razão resolve. --- A Metaphysica, repetimos, não é um tecido de syllogismos forjados \emph{a-priori}; é a investigação das causas \emph{ultimas}, de que dependem os factos. E, como essas causas devem ser e são reaes, porque reaes são os factos, tambem real deve ser e é a Metaphysica, que tem por seu objecto as causas \emph{ultimas} dos entes. --- Concluimos que a existencia de Deus, se é \emph{incognoscivel} pelos sentidos, não é \emph{incognoscivel} pela razão.\par Não menos absurda é a pretenção dos \emph{immanentistas}. Dizem estes --- que Deus não se conhece pela razão, mas pela \emph{consciencia}, --- que é na consciencia que se manifesta o divino, --- que é no intimo da alma que se encontra Deus. --- Mas, como a consciencia é \emph{psychologica} ou \emph{moral}, conforme refere a existencia das nossas operações intellectuaes, ou indica a \emph{conveniencia} d'essas operações com as regras dos costumes, d'ahi a necessidade de perguntarmos aos \emph{immanentistas} se a consciencia, de que fallam, é a \emph{psychologica}, ou é a \emph{moral}? --- Se entendem fallar da consciencia \emph{psychologica}, enganam-se redondamente. Esta consciencia revela-nos as nossas operações, as nossas necessidades e desejos, a existencia do \emph{eu}, a realidade nossa, mas não a realidade de Deus. Diz-nos, é certo, a consciencia que temos a aspiração, o desejo do infinito, de Deus; mas a \emph{aspiração} do infinito não é a \emph{presença} do infinito. Se o infinito se sentisse na nossa consciencia, não haveria aspiração ou desejo; o que se possue, não se deseja. O infinito está fóra e sobre nós; o divino não é \emph{immanente}, é \emph{transcendente}. Fundando-se nestas aspirações ou desejos, a razão, pelo principio de \emph{causalidade}, eleva-se a Deus... Admittimos que Deus está em nós; mas está como \emph{causa} creadora e conservadora do nosso ser; e esta presença de Deus, como \emph{causa} do nosso ser, não se percebe pela consciencia, mas pelo raciocinio, emquanto os seres finitos suppõem e exigem a existencia de um ser infinito... Se teimarmos em procurar a Deus \emph{unicamente} com os dados da consciencia, nunca o encontraremos, a não ser que se admitta --- ou que o ser divino se identifica com o nosso ser, como dizem os \emph{pantheistas}, --- ou que o ser divino, embora distincto do nosso, está sempre presente á nossa intelligencia, como objecto de percepção immediata e directa, como affirmam os \emph{ontologistas}, --- ou que o ser divino é representado e unido á nossa intelligencia por uma idéa ou especie \emph{innata}, como querem os \emph{innatistas}. --- Não menos grosseiro é o engano dos \emph{immanentistas}, se ou quando dizem --- que é a consciencia \emph{moral} que revela o \emph{divino}, --- que Deus vive em nós, como suprema lei moral, --- que a consciencia \emph{moral} é o echo da voz de Deus, é o \emph{proprio Deus que falla}. Por quanto, se sentimos que os dictames da lei moral nos \emph{obrigam}, sentimos tambem que essa \emph{obrigação} não deriva da nossa vontade, mas nos é imposta por um superior, que está sobre e fóra de nós; pois ninguem pode ser subdito e superior de si mesmo. Logo da existencia da lei moral podemos e devemos deduzir a existencia de um legislador supremo... Alem d'isso, a consciencia, sendo a regra \emph{proxima} dos nossos actos, é o \emph{echo} da voz de\pgn{424} Deus, que é a regra \emph{ultima}. E' o echo da voz de Deus; mas não é o proprio Deus; porque o echo de uma voz distingue-se da propria voz, e da pessoa que a emitte. O echo da lei moral só pode servir de base para a conclusão da existencia de um supremo legislador. --- Não podemos tambem negar que a vontade tem o seu logar na demonstração da existencia de Deus. Mas é preciso determinar com exactidão o papel da mesma faculdade nesta demonstração. A vontade pode entrar --- ou para levar a intelligencia a \emph{considerar} os argumentos, que provam a existencia de Deus, --- ou para levar a intelligencia a \emph{admittir cegamente} a existencia de Deus. Na primeira hypothese, o influxo da vontade na intelligencia é apenas \emph{extrinseco} e \emph{mediato}, e isto verifica-se em todos os conhecimentos; na segunda, o influxo é \emph{intrinseco} e \emph{immediato}, e isto verifica-se no acto de Fé, que não é determinado pela \emph{evidencia} da verdade proposta, mas pelo imperio da vontade. Ora é claro que a vontade entra na demonstração da existencia de Deus --- não emquanto leva a intelligencia a \emph{admittir cegamente} essa existencia, mas emquanto a determina a considerar os argumentos, que \emph{provam} tal verdade; de modo que o assenso da intelligencia é devido á força dos argumentos, e não ao imperio da vontade. Esta nunca poderá fazer com que seja verdadeira uma proposição, que a intelligencia julgou falsa. Dizer, pois, com os \emph{immanentistas}, --- que o sentimento deve substituir a razão, --- que a emoção é um conhecimento, --- que a disposição subjectiva deve ter e tem o logar da verdade objectiva, é blasonar um erro tão absurdo quão pueril.\par A Fé confirma as conclusões da razão. O Concilio Vaticano emanou a seguinte definição: « Si quis dixerit \textsc{Deum} unum et verum, Creatorem et Dominum nostrum, \emph{per ea quae facta sunt, naturali rationis lumine certo cognosci non posse}, anathema sit ». --- O Concilio baseou a sua definição no Livro da \emph{Sabedoria} (c. 13, vv. 1-5) e na Epistola de S. Paulo aos Romanos (c. 1, vv. 20-21). --- Estas auctoridades devem ser recebidas, ao menos pelos \emph{tradicionalistas}, que reconhecem a \emph{revelação} como unico criterio de verdade.}.

7. A \textbf{existencia de Deus deve ser demonstrada}. --- A existencia de um ente deve ser demonstrada, quando ella não é percebida \emph{directa} e \emph{immediatamente} --- nem pelas faculdades \emph{sensitivas}, --- \pg{422} \pg{423} nem pela \emph{intelligencia}. Ora a existencia de \textsc{Deus} não é percebida \emph{directa} e \emph{immediatamente} --- nem pelas faculdades \emph{sensitivas}, porque \textsc{Deus} não é um ente sensivel, --- nem pela \emph{intelligencia}, \pg{424} pois, como demonstrámos na \emph{Anthropologia}, a visão \emph{intuitiva} de \textsc{Deus} excede a natural capacidade da nossa intelligencia. Logo a existencia de \textsc{Deus} deve ser demonstrada \footnote{Os que negam a necessidade de demonstrar a existencia de \textsc{Deus} são --- \emph{a}) os \emph{sentimentalistas}, sustentando que conhecemos \emph{immediatamente} a \textsc{Deus} por meio de um sentido arcano e divino. --- \emph{b}) os \emph{ontologistas}, admittindo uma \emph{intuição immediata} ou \emph{directa} de \textsc{Deus}.\par Estas duas opiniões são falsas.\par \emph{a}) E' falsa a dos \emph{sentimentalistas}; pois a existencia d'esse sentido arcano e divino é contraria ao testemunho da consciencia.\par \emph{b}) E' falsa a dos \emph{ontologistas}; pois o \emph{ontologismo}, como vimos, é falso.\par Aqui apenas accrescentaremos que a S. Egreja, Mestra infallivel de verdade, declarou suspeitas e perigosas, entre outras, estas proposições: « Immediata \textsc{Dei} cognitio, habitualis saltem, intellectui humano essentialis est, ita ut sine ea nihil cognoscere possit: siquidem est ipsum lumen intellectuale ». --- « Esse illud, quod in omnibus et sine quo nihil intelligimus, est esse divinum ». (Decr. S. Officii, die 18 Septemb. a. 1861).}.

8. \textbf{A existencia de Deus não se demonstra a-priori.} --- A demonstração diz-se \emph{a-priori}, quando a existencia de um ser se prova\pg{425} pela sua \emph{causa efficiente}. Ora a existencia de \textsc{Deus} não se prova pela sua \emph{causa efficiente}; porque \textsc{Deus} é a Causa primeira, e por isso absolutamente independente de toda e qualquer causa. Logo a existencia de \textsc{Deus} não se demonstra \emph{a-priori}.

9. \textbf{A existencia de Deus não se demonstra a-simultaneo.} --- A demonstração é \emph{a-simultaneo}, quando a existencia de um ente se prova pela analyse da idea ou noção do proprio ente. Ora este processo, applicado á existencia de \textsc{Deus}, é vicioso; visto que, para uma coisa existir na realidade, não basta que exista no pensamento, pois deduzir da ordem \emph{ideal} a ordem \emph{real} é tirar uma consequencia mais extensa que as premissas. Logo a existencia de \textsc{Deus} não se demonstra \emph{a-simultaneo} \footnote{Esta proposição --- \textsc{Deus} \emph{existe} --- é \emph{immediatamente} evidente \emph{em si mesma}, mas não é \emph{immediatamente} evidente \emph{em relação a nós}.\par a) E' \emph{immediatamente} evidente \emph{em si mesma}. Uma proposição é \emph{immediatamente} evidente \emph{em si mesma}, quando o predicado convém \emph{essencialmente} ao sujeito. Ora a \emph{existencia} convém \emph{essencialmente} a \textsc{Deus}; porque \textsc{Deus} é o ser \emph{necessario}, e não póde deixar de existir. Logo a proposição --- \textsc{Deus} \emph{existe} --- é \emph{immediatamente} evidente em si mesma.\par \emph{b}) Não é \emph{immediatamente} evidente \emph{em relação a nós}. Uma proposição é \emph{immediatamente} evidente em relação a nós, quando directamente e sem necessidade de raciocinio conhecemos a essencia do sujeito e do predicado; assim a proposição --- \emph{o todo é maior que a parte} --- é \emph{immediatamente} evidente \emph{em relação a nós}, pois basta enunciar as palavras \emph{todo} e \emph{parte}, para logo percebermos que o \emph{todo} é maior que a \emph{parte}. Ora a essencia de \textsc{Deus}, não sendo por nós conhecida senão depois de um raciocinio, e ainda assim muito imperfeitamente, não se nos manifesta logo á primeira vista como exigindo a existencia. Logo a proposição --- \textsc{Deus} \emph{existe} --- não é \emph{immediatamente} evidente \emph{em relação a nós}. --- Cf. S. Thom., --- \emph{Sum. Th.}, p. I, q. 2, a. 1; --- \emph{c. Gentes.}, l. I, cc. 10, 11; 1 \emph{Sent.} Dist. 3, q. 1, a. 2; \emph{De verit.}, q. 10, a. 12.\par Os mais celebres argumentos \emph{a-simultaneo} em favor da \emph{existencia} de \textsc{Deus} são os de \emph{Santo Anselmo}, \emph{Descartes} e \emph{Leibnitz}.\par a) \emph{Santo Anselmo} raciocina d'este modo: Uma proposição é \emph{immediatamente evidente em relação a nós}, quando a sua verdade se nos manifesta pela simples enunciação dos termos. Ora a verdade d'esta proposição --- \textsc{Deus} \emph{existe} --- torna-se-nos manifesta pela simples enunciação dos termos. Com effeito, quando se profere o nome de \textsc{Deus}, immediatamente entendemos que \textsc{Deus} é o Ser maior que se póde imaginar, e que por isso deve existir na \emph{realidade}; porque, se existisse unicamente na intelligencia, não seria o Ser maior que se póde imaginar, pois superior a Elle seria o ser, que existisse na intelligencia e tambem na \emph{realidade}.\par Critica. Ainda que todos admittissem que \textsc{Deus} é o ser maior que se póde imaginar, todavia d'ahi só poderia concluir-se que \textsc{Deus} existe na \emph{intelligencia}, mas não na \emph{realidade}; porque, como dissemos, deduzir da ordem \emph{ideal} a ordem \emph{real} é tirar uma consequencia, que não está incluida nas premissas\pgn{426}. --- Para provar que \textsc{Deus} existe tambem na \emph{realidade}, deve demonstrar-se que existe na \emph{realidade} um ser, que é o maior que se póde imaginar. --- Cf. S. Thom., \emph{c. Gentes}, l. I, c. 11.\par b) \emph{Descartes} diz: Todo o attribuito, que está contido na idea \emph{clara} e \emph{distincta} de um sujeito, deve affirmar-se do mesmo sujeito. Ora na idea \emph{clara} e \emph{distincta} de \textsc{Deus} está contida a existencia \emph{real}; porque \textsc{Deus} é infinitamente perfeito, e a existencia \emph{real} é uma perfeição. Logo deve affirmar-se de \textsc{Deus} a existencia \emph{real}, isto é, \textsc{Deus} existe \emph{realmente}.\par Critica. O argumento de \emph{Descartes} tem o mesmo defeito que o de \emph{Santo Anselmo}; pois passa da ordem \emph{ideal} para a \emph{real}. Na verdade, tudo o que está contido na idea \emph{clara} e \emph{distincta} de um sujeito deve attribuir-se ao sujeito \emph{na mesma ordem}, em que está o sujeito; de modo que, se o sujeito estiver na ordem \emph{real}, o predicado tambem estará na ordem \emph{real}, --- mas, se o sujeito estiver na ordem \emph{ideal}, o predicado tambem existirá necessariamente na ordem \emph{ideal}. Logo na \emph{idea} de \textsc{Deus} está contida a existencia \emph{ideal}, e não a \emph{real}. --- Lembramos que, para \emph{Descartes}, a idea do \emph{infinito} é \emph{innata} (o que é falso), e que a verdade da proposição maior do seu syllogismo se basêa na veracidade do proprio \textsc{Deus} (o que constitue um \emph{circulo vicioso}).\par c) \emph{Leibnitz} apresenta este argumento: O ente infinitamente perfeito é possivel. Se é possivel, existe na \emph{realidade}; porque a existencia \emph{real}, por ser uma perfeição, pertence necessariamente á essencia do Ente infinitamente perfeito. Logo o Ente infinitamente perfeito existe na \emph{realidade}.\par Critica. Este argumento é inconcludente pelo mesmo defeito, que apontámos nos dois antecedentes. Do facto de \emph{concebermos}, como \emph{possivel}, o Ente infinitamente perfeito, cuja essencia exige a existencia, só resulta que não podemos \emph{conceber}, como \emph{possivel}, esse ente, sem que o \emph{concebamos} como existente; mas não se segue que a realidade deva corresponder ao que \emph{concebemos} na intelligencia.}.

10. \textbf{A existencia de Deus demonstra-se a-posteriori.} --- A demonstração \emph{a-posteriori} faz-se, quando, pela existencia do effeito,\pg{426} se prova a existencia da causa respectiva. Ora a existencia de \textsc{Deus} demonstra-se pela existencia das creaturas; porque estas, por serem \emph{effeitos}, suppõem uma causa proporcionada, que é \textsc{Deus}. Logo a existencia de \textsc{Deus} demonstra-se \emph{a-posteriori} \footnote{Cf. S. Thom., --- \emph{Sum}. \emph{Th}., p. I, q. 2, a. 2; --- \emph{c}. \emph{Gentes}, l. I, c. 12.\par A demonstração da existencia de Deus basea-se no \emph{principio de causalidade}. E' por isso que os \emph{materialistas}, os \emph{positivistas} e todos os \emph{incredulos}, no intuito de negar esta verdade fundamental, ou, ao menos, de prescindir d'ella, atacam a noção de \emph{causa} e esse principio de \emph{causalidade}. --- Não repetimos tudo o que na \emph{Ontologia} escrevemos em defeza do dito principio. Limitamo-nos a observar que o principio de \emph{causalidade} é uma verdade, que a intelligencia descobre e formúla á vista das mudanças, dos phenomenos, dos factos, que se succedem, uns aos outros, no mundo visivel. Ora quando um phenomeno começa, percebe-se immediatamente que está relacionado com um outro, de que deriva. Esta relação não é apenas de \emph{successão}, mas de \emph{dependencia}. Tal dependencia é uma coisa real, objectiva, existente nos proprios seres creados; a intelligencia \emph{descobre-a}, mas \emph{não a inventa}. Por isso o principio de \emph{causalidade} tem um valor real e objectivo; porque, embora este juizo --- \emph{todo o effeito tem a sua causa} --- seja um producto da razão, comtudo a \emph{razão não o produz senão fundada nos dados da experiencia sensivel}. A experiencia apresenta os factos; a razão descobre o nexo, que os liga. ---\pgn{427}Advirta-se que o principio de \emph{causalidade} não se applica a todo o ser que \emph{existe}, mas a todo o ser que \emph{começa} a existir, e por isso a Causa \emph{primeira} não tem causa.\par Os \emph{Kantistas} não negam o principio de \emph{causalidade}, mas dizem que elle só tem valor, quando se trata de dois phenomenos, ambos sujeitos á observação sensivel, porque, em tal caso, podemos averiguar a dependencia de um phenomeno do outro, --- mas não tem valor, quando se trata de duas realidades, uma das quaes não está sujeita á experiencia, porque, em tal caso, não podemos conhecer a \emph{dependencia} ou a \emph{origem} de uma realidade da outra, e por isso nem a razão de \emph{causa} e de \emph{effeito}; o que se verifica no nosso caso, porque, se o mundo está sujeito á experiencia, Deus escapa a toda a experiencia. Por isso não podemos basear-nos no principio de \emph{causalidade}, para deduzir da existencia do mundo a existencia de Deus.\par Respondemos que, para o principio de \emph{causalidade} ter todo o seu valor, não é necessario que ambas as realidades, que se dizem proceder uma da outra, estejam sujeitas á experiencia sensivel. Basta que um phenomeno, um facto, presente aos nossos sentidos, manifeste um caracter de mutabilidade e de contingencia, para que possamos e devamos concluir que esse phenomeno \emph{depende} necessariamente de um antecedente. Pouco importa, para o caso, que esse antecedente seja visivel ou invisivel, natural ou sobrenatural; elle existe, deve existir, porque existe o phenomeno, que não tem em si a razão sufficiente da sua existencia. Ninguem viu a força da \emph{gravitação}; e, comtudo, todos a admittem, porque se revela na queda dos graves. Ninguem viu o talento militar de Napoleão; mas todos o admittem, porque se manifestou nos seus grandes planos de guerra. Deus é invisivel na sua essencia, mas revela-se nas coisas creadas, que suppõem e exigem uma força infinita. Portanto, ou devemos negar o principio de \emph{causalidade}, e temos a completa derrocada de todas as sciencias, --- ou devemos admittir que, em força d'esse principio, podemos e devemos elevar-nos ao conhecimento da existencia de Deus. --- Nem se diga que uma percepção mais profunda e completa dos phenomenos naturaes pode fazer com que se prescinda de uma causa primeira. Os phenomenos naturaes, ainda que podessem ser percebidos em toda a sua comprehensão, nunca podem perder os \emph{essenciaes} caracteres de \emph{contingencia} e de \emph{mutabilidade}, e por isso da sua \emph{dependencia} de uma causa primeira. Pelo contrario: quanto mais se conhece o \emph{movimento}, a \emph{harmonia}, a \emph{finalidade} dos seres visiveis, tanto mais necessaria apparece a existencia de um supremo e omnipotente Ordenador.\par Objectam: Todos os nossos conhecimentos são \emph{relativos} ao mundo \emph{visivel} e á intelligencia \emph{finita} do homem. Logo não podemos conhecer a \textsc{Deus}, que se suppõe \emph{invisivel} e \emph{infinito}.\par Resposta. E' certo que todos os nossos conhecimentos, \emph{que adquirimos}\pgn{428}\emph{de um modo directo e immediato}, se referem ao mundo visivel. Mas d'ahi não se segue que a nossa razão não possa, nem de um modo \emph{indirecto} e \emph{mediato}, adquirir outros conhecimentos, relativos ao mundo invisivel. Os caracteres de contingencia e de mutabilidade, que observamos no mundo visivel, levam-nos necessariamente a admittir a existencia de um Ente invisivel, necessario e immutavel. --- Como tambem é certo que todos os conhecimentos são proporcionados á capacidade da intelligencia cognoscente e ao modo por que os objectos se tornam presentes á mesma intelligencia, de modo que só apprehendemos as notas do objecto, que actuaram na faculdade, ficando incognosciveis as que escaparam ao alcance d'ella. Mas esta \emph{relatividade}, se prova que não podemos \emph{comprehender} a Deus, porque a intelligibilidade de Deus excede infinitamente a capacidade da intelligencia creada, não prova que não possamos, \emph{de algum modo}, conhecer o omnipotente Creador de todas as coisas.\par Devemos, porém, advertir que, para os sequazes do \emph{agnosticismo}, \emph{relatividade} significa \emph{subjectivismo}; e por isso quando dizem que os nossos conhecimentos são \emph{relativos}, entendem que são \emph{subjectivos}, em quanto nos fazem conhecer as coisas não como são em si mesmas, mas como são no nosso espirito, ás quaes nada corresponde na realidade. Não nos detemos na refutação d'este erro, a qual já está feita. Só accrescentaremos que, se todos os nossos conhecimentos são \emph{subjectivos}, são \emph{subjectivos} não só os que se referem ao mundo invisivel, mas tambem os que se referem ao visivel, e por isso, se não existe Deus, nem existe coisa alguma no mundo. --- Nem se diga que os conhecimentos relativos ao mundo \emph{visivel} são \emph{objectivos}, porque podem ser verificados pela experiencia, pelos factos: o que se não pode fazer com os conhecimentos relativos ao mundo \emph{invisivel}, a Deus. Negamos que os nossos conhecimentos relativos a Deus não se possam verificar pelos factos. Na verdade, como provamos nós que Deus existe? Provamol-o pela origem e natureza do mundo, do movimento, da vida, da ordem, da finalidade; porque a origem e a natureza d'estas coisas suppõem a existencia de uma causa extramundana. Ora o mundo, o movimento, a vida, a ordem, a finalidade, são factos, e factos grandiosos, admiraveis, que actuam nos nossos sentidos e despertam a admiração na nossa intelligencia. Portanto o conhecimento de Deus é tão \emph{objectivo}, como \emph{objectivos} são os factos, que nos servem de degrau para nos elevarmos a Elle.\par Continuam: A existencia de Deus não pode ser demonstrada pelos effeitos. Na verdade, entre a causa e o effeito deve existir uma proporção; ora nenhuma proporção pode existir entre as creaturas, que são seres finitos, e Deus, que é o Ser infinito.\par Resposta. Demonstra-se a causa pelos effeitos por dois modos. Se os effeitos possuem a mesma natureza da causa (o que se verifica nas causas \emph{univocas}), então não só levam ao conhecimento da \emph{existencia}, mas tambem da \emph{natureza} da causa, de modo que adquirimos um conhecimento \emph{perfeito} da causa. Se os effeitos não têm a mesma natureza da causa (o que se verifica nas causas \emph{equivocas}), então só levam ao conhecimento da \emph{existencia} da causa, e não da \emph{natureza}, e por isso o conhecimento da causa é \emph{imperfeito}\pgn{429}, e o grau d'esta imperfeição é proporcionado ao grau de dessemelhança entre os effeitos e a causa. Por isso, \emph{todo} o effeito tem com a sua causa uma proporção de \emph{ordem}, emquanto é produzido por ella; mas nem \emph{todo} o effeito tem com a causa uma proporção de \emph{natureza}. Appliquemos estes principios ao nosso caso. Os seres visiveis, sendo finitos e contingentes, não exprimem toda a virtude da sua causa, não possuem a mesma \emph{natureza} de Deus, e por isso não nos levam ao conhecimento \emph{adequado} da \emph{natureza} de Deus. Mas, sendo \emph{effeitos}, e effeitos que só podem derivar de um Ente necessario e omnipotente, levam certamente ao conhecimento da \emph{existencia} d'esse Ente. --- Todavia, embora o conhecimento, que temos da \emph{natureza} de \textsc{Deus}, não seja \emph{adequado}, mas \emph{imperfeito}, comtudo é \emph{positivo}; pois as coisas creadas são semelhanças \emph{positivas}, embora \emph{imperfeitas}, das divinas perfeições, visto que todo o agente tende a imprimir no effeito a propria \emph{semelhança}.}.

%%>>> art-37

\section*{Artigo II.\\ Demonstração da existencia de Deus.}
\addcontentsline{toc}{section}{Artigo II. --- Demonstração da existencia de Deus.}

11. \textbf{Especies de argumentos para a demonstração da existencia de Deus}. --- A existencia de \textsc{Deus} demonstra-se \emph{a-posteriori}. Ora os argumentos, que demonstram \emph{a-posteriori} a existencia de \pg{427} \pg{428} \textsc{Deus}, são de tres especies: \emph{metaphysicos}, \emph{physicos} e \emph{moraes}. Os \emph{metaphysicos} basêam-se na propria \emph{natureza} das coisas, --- os \emph{physicos}, na \emph{ordem} que existe no mundo e, em particular, nas \pg{429} \emph{operações} dos agentes \emph{irracionaes}, --- os \emph{moraes}, na existencia da lei \emph{moral} e no \emph{consenso} de todos os povos \footnote{Na classificação dos argumentos, que demonstram a existencia de \textsc{Deus}, em \emph{methaphysicos}, \emph{physicos}, \emph{e moraes}, conforme se baseam na \emph{essencia} das coisas, ou na \emph{ordem} existente no mundo e nas \emph{operações} dos agentes \emph{irracionaes}, ou na lei \emph{moral} e no \emph{consenso} dos povos, seguimos a doutrina dos mais sensatos auctores.\par Alguns dividem esses argumentos em \emph{cosmologicos}, \emph{psychologicos} e \emph{moraes}; conforme os factos, em que se fundam, são \emph{cosmologicos} (baseados nas leis \emph{physicas}), \emph{psychologicos} (baseados nas leis \emph{intellectuaes}), ou \emph{moraes} (baseados nas leis \emph{moraes}). --- Esta divisão tambem é acceitavel.\par \emph{Kant}, dizendo que o espirito humano se eleva a \textsc{Deus} por meio de um d'estes tres principios: de \emph{identidade}, de \emph{causalidade} ou de \emph{finalidade}, reduz as provas a tres classes: \emph{ontologicas}, \emph{cosmologicas} e \emph{physico-teleologicas}, conforme se baseam no principio de \emph{identidade}, de \emph{causalidade} ou de \emph{finalidade}. --- A doutrina de \emph{Kant} é inacceitavel. O principio, em que se baseam todas as provas em favor da existencia de \textsc{Deus}, é um só: o \emph{principio de causalidade}, e a distincção das provas deriva da natureza dos \emph{factos}, aos quaes o dito principio se applica. O principio de \emph{identidade} não prova coisa alguma, como o proprio \emph{Kant} confessa; o principio de \emph{finalidade}, longe de ser a base do argumento \emph{teleologico}, é apenas um corollario d'esse argumento.\par A existencia de \textsc{Deus} é uma verdade tão necessaria e fundamental, que recorre em todos os grandes problemas da philosophia. Nem admira. Não encontrando nos agentes finitos a causa \emph{sufficiente} dos seres, a razão deve elevar-se a uma causa \emph{suprema} e \emph{infinita}. A \emph{origem} do \emph{mundo} (\emph{Cosm}., \emph{ger}.\pgn{430}Cap. I, art. V), a \emph{finalidade} nas operações dos seres irracionaes (\emph{Cosm}. \emph{ger}., Cap. II, art. III), a ordem \emph{physica} com as suas \emph{leis} (\emph{Cosm}. \emph{ger}., Cap. II, art. IV e V), a \emph{origem} e o \emph{desenvolvimento da vida} (\emph{Cosm}. \emph{esp}., Cap. II, art. III), a \emph{origem} da \emph{alma} e do \emph{corpo} humano (\emph{Anthr}., Sec. IV, Cap. I), tudo o que nos cerca, com os seus caracteres de \emph{contingencia} e de \emph{mutabilidade}, proclama a existencia de uma Causa infinitamente sabia, poderosa e boa --- a existencia de \textsc{Deus}.\par Escutámos esta voz maravilhosa da creação e podemos dizer, com immensa satisfação da nossa alma, que procurámos sempre pôr em relevo a racionabilidade e a necessidade da existencia de um Ente supremo, e mostrar a fatuidade, os erros e as contradicções dos adversarios.\par Mas isso não obsta a que façamos aqui, como no logar proprio, uma synthese do que dissemos em todo este trabalho, indicando os principaes argumentos em defeza da \emph{existencia} de Deus. --- Cf. \emph{Sum}. \emph{Th}., p. I, q. 2, a. 3; --- \emph{c}. \emph{Gentes}, l. I, c. 13.}.

12. \textbf{Argumentos metaphysicos.} --- Os argumentos \emph{metaphysicos}, em favor da \emph{existencia} de \textsc{Deus}, baseam-se em que, --- \emph{a}) se no mundo existe o movimento ou mudança, deve existir um motor \emph{immovel}, --- \emph{b}) se existem as causas \emph{segundas}, deve existir a causa \emph{primeira}, --- \emph{c}) se existem seres \emph{contingentes}, deve existir um ser \emph{necessario}, --- \emph{d}) se existem varios e diversos graus de perfeição, deve existir um ser infinitamente perfeito, --- \emph{e}) se existem verdades \emph{absolutas} e \emph{necessarias}, deve existir uma intelligencia \emph{absoluta} e \emph{necessaria}.
\pg{430} a) As substancias creadas estão sujeitas a continuos movimentos ou mudanças. Ora os movimentos, sendo passagens da \emph{potencia} para o \emph{acto}, devem ser determinados --- não pelas proprias substancias que se movem, pois ninguem dá o que não tem, --- mas por um motor distincto, pois o ente, dotado apenas de \emph{potencia}, não adquire o \emph{acto}, se lh'o não communica um outro ente, que já o possue. Este motor --- ou é \emph{movido}, --- ou é \emph{immovel}. Se é \emph{immovel}, é \textsc{Deus}. Se é \emph{movido}, subiremos de motor em motor, até encontrarmos necessariamente um motor \emph{immovel}; porque, se não existisse este motor \emph{immovel}, não poderiam existir os motores \emph{moveis}, que não têm em si a razão do seu movimento, e por isso não haveria movimento nas substancias creadas. Logo existe \textsc{Deus} \footnote{Aqui a palavra \emph{movimento} toma-se em sentido mais lato, emquanto exprime não só o \emph{movimento local}, mas toda e qualquer \emph{mudança} ou \emph{passagem da potencia para o acto}. D'este modo o argumento exposto equivale a este: O mundo é \emph{mutavel}. Ora todo o ente \emph{mutavel} depende do ente \emph{immutavel}. Logo existe o ente \emph{immutavel}. O ente \emph{immutavel} é \textsc{Deus}. Logo existe \textsc{Deus}. --- Ora, encontrando-se a \emph{mudança} não só nos corpos mas tambem nos espiritos, o argumento póde desdobrar-se; pois, baseando-nos tanto no movimento dos \emph{corpos} como no dos \emph{espiritos}, podemos elevar-nos ao descobrimento de uma causa, que move tudo, sendo ella \emph{immovel} ou \emph{immutavel}.\par Eis os dois argumentos:\par a) A experiencia attesta que as substancias materiaes, por exemplo, os planetas, estão sujeitas a continuos movimentos. Ora todo o corpo que se move, é movido por um agente; pois o corpo, sendo por si indifferente tanto para a quietação como para o movimento, não póde determinar-se por si ao movimento. Logo ou devemos admittir uma serie infinita de motores \emph{moveis}, ou deve existir um motor \emph{immovel}, que seja a causa de todo o movimento dos corpos. Ora uma serie infinita de motores \emph{moveis} repugna; porque numa serie \emph{infinita} de motores \emph{moveis} não haveria um \emph{primeiro} motor, e assim não existiria uma causa proporcionada ou adequada de todos os movimentos\pgn{431}nos corpos. Logo deve existir um motor \emph{immovel}. Este motor \emph{immovel} deve ser infinitamente perfeito; pois, se o não fosse, se lhe faltasse alguma perfeição, poderia tender para ella, e assim passaria da potencia ao acto, e não seria \emph{immovel}. O ser infinitamente perfeito é \textsc{Deus}. Logo existe \textsc{Deus}.\par Este argumento, muito popular, pode desenvolver-se \emph{scientificamente} da maneira seguinte. Existe no mundo o movimento. Para a explicação da existencia do movimento adduzem-se duas hypotheses: a \emph{mechanica} e a \emph{dynamica}. Ora nenhuma d'estas hypotheses pode, na explicação do movimento, prescindir da intervenção de um primeiro Motor. --- Não pode prescindir a hypothese \emph{mechanica}. Na verdade, para os sequazes d'esta hypothese, a materia é \emph{meramente passiva}, de modo que os corpos são movidos por um impulso extrinseco, mas elles não movem, e só communicam aos outros corpos o movimento recebido. Se todos os corpos, de que se compõe o universo, são meramente passivos, é necessario que a força impulsiva do movimento venha de fóra, e proceda de um primeiro Motor, que não seja movido. Nem se diga que o movimento da materia teve um começo \emph{necessario}, ou que é \emph{eterno}; porque um \emph{começo necessario} é uma contradicção (o que começa é necessariamente contingente), e um movimento \emph{eterno} da materia, embora não envolvesse repugnancia, não deixaria de exigir um primeiro motor, desde que se supponha a materia meramente \emph{passiva}. --- Não pode prescindir a hypothese \emph{dynamica}. Com effeito, para os sequazes d'esta hypothese, a materia não é meramente passiva, mas é essencialmente dotada de actividade, em virtude da qual um corpo pode, em certas circumstancias, mover outro corpo, \emph{attrahendo-o}. E' certo que o mundo inteiro se move pela \emph{lei da attracção}; mas a \emph{attracção} não é o \emph{primeiro} motor, nem é o motor \emph{sufficiente} do nosso systema solar. Não é o \emph{primeiro} motor; porquanto, a força de \emph{attracção}, de que é dotada a materia, não está sempre e necessariamente \emph{em acto}, mas, pelo contrario, tende para o repouso das partes numa unica massa, e por isso é indifferente para o acto e para a potencia, e, não conténdo em si a razão da passagem da potencia para o acto, exige e suppõe uma causa que a determine e mova. Não é o motor \emph{sufficiente} do nosso systema solar; porquanto este systema deve ser mantido em equilibrio e regulado nos seus movimentos por duas forças --- uma \emph{centripeta} (sem a qual não haveria unidade no systema) e outra \emph{centrifuga} (sem a qual todos os corpos confundir-se-hiam numa unica massa), e se a força \emph{centripeta}, da qual deriva a attracção, pode dizer-se essencial á materia, não se pode dizer essencial a força \emph{centrifuga}, a qual por isso deve ser accrescentada e unida por um motor extranho e superior á propria materia.\par Dizem que este argumento, se prova a existencia de um Ente superior ao mundo e causa de todo o movimento, não prova a existencia de um Deus tão perfeito, como o concebem os espiritualistas. --- Respondemos que o nosso argumento, embora não provasse senão a \emph{existencia} de um Ser supremo, já teria alcançado o seu fim, que é o de demonstrar a mesma \emph{existencia}, e não a \emph{natureza} do Ente supremo. São duas questões distinctas. --- Todavia devemos admittir que este argumento, que prova a necessidade de um \emph{primeiro Motor}, contém, \emph{em germe}, a razão de todos os attributos da Divindade.\pgn{432}Na verdade, se a materia, por ser indifferente ao repouso e ao movimento, tem a razão do seu movimento não em si mas numa causa exterior, segue-se que tal causa existe fóra da materia, e por isso é immaterial, intelligente, eterna, necessaria, e tambem infinita e unica. Porquanto, o primeiro Motor, não podendo ser movido, é \emph{todo acto}, sem mistura de potencialidade, é acto \emph{purissimo}, isto é, infinitamente perfeito e é unico.\par Os melhores naturalistas e philosophos modernos reconhecem esta verdade. Citaremos apenas \emph{Couchy}, \emph{Clarke}, \emph{B}. \emph{S}. \emph{Hilaire}, \emph{P}. \emph{Janet}. No fim, ouviremos um insuspeito, \emph{Rousseau}.\par \emph{Couchy} diz: « O unico ser, que póde produzir as forças physicas, é o Ente necessario. A força é a expressão da sua vontade. As diversas forças do equilibrio e movimento são causas secundarias. A gravitação universal, a attracção dos corpos na razão inversa do quadrato das distancias, que os separam, é uma lei que exprime a vontade divina » (\emph{Comptes rendus hebd}. \emph{de l'Acad}. \emph{des sciences}, t. XXI, Paris, 1845).\par \emph{Clarke}, discipulo de \emph{Newton}, fallando da \emph{attracção}, de que depende todo o movimento dos corpos celestes, escreve: « Ella póde ser o effeito de um impulso, mas não póde ser material, é antes o effeito de uma causa immaterial » (\emph{Physica de Rohault}, trad. por \emph{Clarke}, II, c. XI).\par \emph{Barthélemy Saint-Hilaire} diz: « Concebe-se que a materia siga hoje as leis infalliveis, que a regem, e conserve uma regularidade eterna; mas foi necessario um impulso primitivo, que ordenou tudo para a inexhaurivel serie dos tempos. Até as anomalias, que offerece o admiravel systema dos ceus, attestam a presença immanente e indefectivel d'Aquelle que os fez. O sobrenatural está em tudo... Só é necessario que a sciencia se resigne a resolver certos problemas por outros meios, que não sejam os de uma experiencia impossivel; e o da origem de todas as coisas é um d'esses problemas, a que só se renuncia por timidez, pensando que se pratica uma reserva prudente. A \emph{questão da origem é inevitavel}, e de nada valeria querer illudil-a ». (\emph{Journal des Savants}, 1862, pag. 607).\par \emph{P}. \emph{Janet} diz: « A experiencia não mostra senão que um corpo é movido por outro, e este por outro, e isto ao infinito. Ora, se fóra d'esta cadêa não existe uma causa motora, teremos uma serie de movimentos sem causa e sem razão » (\emph{Le matérialisme contemporain}, 1875, pag. 58).\par \emph{Rousseau} escreve: « Vejo a materia ora em movimento e ora em repouso; d'ahi concluo que nem o repouso nem o movimento lhe são essenciaes. Mas o movimento, sendo uma acção, mostra ser effeito de uma causa... Conceber a materia productora do movimento é conceber um effeito destituido de causa... Não é evidente que, se o movimento fosse essencial á materia, um não poderia separar-se da outra, mas estaria sempre unido com ella, conservando-se sempre no mesmo grau, sempre o mesmo em cada uma das partes, não poderia augmentar nem diminuir, e não se poderia conceber a materia em repouso?!... Quando me dizem que o movimento não é \emph{essencial} á materia, mas \emph{necessario}, pretendem embaraçar-me com palavras, que poderiam facilmente confutar-se, se tivessem algum sentido. Porque o movimento da materia --- ou deriva da propria materia, e então lhe é essencial\pgn{433} --- ou deriva de uma causa extrinseca, e já não é necessario á propria materia, senão emquanto a causa motora actua nella » (\emph{Emilio}, t. III).\par Na \emph{Cosmologia geral} tambem desenvolvemos este ponto.\par b) No mundo não só existe o movimento \emph{local}, mas tambem o movimento \emph{espiritual}; e este, como aquelle, suppõe e exige necessariamente a existencia de um primeiro motor immovel. Na verdade, se reflectirmos em nós mesmos, percebemos que a nossa alma tambem se move, emquanto passa continuamente da potencia para o acto. O acto é um complemento, uma addição que aperfeiçôa a potencia. Ora este acto não pode ser e não é extrahido do fundo da potencia, porque ninguem dá o que não tem; mas deve provir de um agente extrinseco, de um motor, que possua \emph{actualmente} o que falta á potencia e que seja capaz de lh'o communicar. Este motor é a propria alma, que se move a si mesma emquanto uma faculdade, que está \emph{em acto}, move uma outra, que está \emph{em potencia}. Todavia a alma não pode ser o \emph{primeiro} motor. O primeiro acto da faculdade, que move todas as outras, não pode ser o effeito do impulso das outras faculdades, que estão em potencia, mas deve ser o effeito de um motor distincto e superior á propria faculdade. Este motor --- ou é o primeiro Motor, --- ou é um motor intermedio, movido, finalmente, por esse primeiro Motor immovel.\par Este argumento, que se refere ao movimento da alma humana em geral, conserva toda a sua força, quando se refere ao movimento da \emph{intelligencia} e da \emph{vontade}.\par a) Dá-se o movimento na \emph{intelligencia}. Ainda que a intelligencia mova a vontade, e esta mova aquella, comtudo, por ser absurdo o processo ao infinito, devemos admittir que o primeiro acto da intelligencia precedeu o primeiro acto da vontade. Diz S. Thomaz: « Non oportet procedere in infinitum, sed statur in intellectu, sicut in primo. Omnem enim voluntatis motum necesse est quod praecedat apprehensio; sed non omnem apprehensionem praecedit motus voluntatis » (\emph{Sum}. \emph{Th}., p. I, q. 82, a. 4, ad 3). Mas de que modo passa a intelligencia da potencia para o acto? Espontaneamente? Não; porque faltaria a razão sufficiente da passagem. Pela acção do objecto intelligivel? Tambem não; porque, sendo este objecto formado pela intelligencia agente, a acção do objecto é \emph{posterior} á da intelligencia. A causa motora deve ser superior ao objecto e a sujeito, ao intelligivel e ao intelligente; e esta causa é Deus. Aristoteles, ainda que pagão, exprime esta verdade nas seguintes magnificas palavras: « Quod est principium motus in anima? Patet vero, quemadmodum in universo est Deus, sic etiam in anima; moventur enim omnia ab inexistente in nobis Numine. Intellectus autem principium non intellectus sed praestantius quidpiam est. Quid autem et scientia praestantius est et ratione, nisi Deus? (\emph{Moral}. \emph{Eudem}., l. 7, c. 14). D'onde se vê a verdade do aphorismo: \emph{Eu penso}; \emph{logo Deus existe}.\par \emph{b}) Não só o acto da intelligencia, mas tambem o da \emph{vontade} exige e suppõe a existencia de um primeiro Motor. A vontade é movida necessariamente pelo bem absoluto. Esse movimento provoca uma tendencia \emph{universal}, \emph{necessaria}, \emph{indestructivel}. Ora um effeito, dotado de semelhantes qualidades, como é essa tendencia, não pode derivar de uma causa imaginaria ou phantastica, mas deve derivar de uma causa real e objectiva, que deve\pgn{434}identificar-se com o proprio bem absoluto; porque todo o effeito suppõe uma causa proporcionada. Ora o bem absoluto é Deus. Logo Deus existe.\par Como se vê, o principio --- \emph{tudo o que se move é movido por um motor distincto e}, \emph{em ultima analyse}, \emph{immovel} --- é muito fecundo e serve para a intelligencia e resolução de innumeraveis problemas.}.

b) No mundo existem muitas e diversas \emph{causas efficientes}. Estas causas estão por tal modo subordinadas entre si, que formam uma serie, na qual uma depende de outra, a inferior \pg{431} \pg{432} é \emph{effeito} da superior. Ora, nesta serie de causas \emph{efficientes}, deve existir uma causa \emph{primeira}, que seja só \emph{causa} e não \emph{effeito}, que seja causa de todas as causas; visto que, se não existisse uma \pg{433} causa \emph{primeira}, não poderiam existir, como existem, as causas \emph{segundas}, pois as causas \emph{segundas} são \emph{effeitos}, e todo o effeito \pg{434} exige necessariamente uma causa proporcionada. A causa \emph{primeira} é \textsc{Deus}. Logo existe \textsc{Deus}\footnote{Os \emph{materialistas}, para destruirem a força d'este argumento, recorrem a uma serie \emph{infinita} de causas \emph{efficientes}, das quaes uma produza outra, sem sermos obrigados a admittir uma \emph{primeira} causa \emph{efficiente}.\par O recurso é vão. Porquanto, essa serie infinita, embora fosse \emph{possivel}. estaria disposta ou em linha \emph{recta}, ou em linha \emph{circular}. Ambas as hypotheses são absurdas. --- E' absurda a \emph{primeira;} porque, se estivesse disposta em linha \emph{recta}, de modo que a causa inferior fosse produzida pela superior, essa serie não teria uma causa \emph{primeira} (porque o infinito não tem principio). Não tendo uma causa \emph{primeira}, não teria a \emph{segunda}, nem a \emph{terceira}, nem causa alguma; visto que numa serie de causas, dispostas em linha \emph{recta}, a inferior não póde existir sem a superior. Logo deve admittir-se, fóra d'esta serie, uma causa primeira. --- E' absurda a segunda hypothese. Na verdade, se a serie \emph{infinita} das causas \emph{efficientes} estivesse disposta em linha \emph{circular}, um e o mesmo ente seria causa e effeito de outro ente. Na verdade, supponhamos uma linha \emph{circular}, formada por a, b, c, d... E' evidente que, sendo a causa de b, b de c, c de d, d de a, d seria causa e effeito (ainda que remoto) de a; o que é absurdo.\par Os adversarios dizem que a \emph{primeira} causa \emph{efficiente} é a \emph{materia}.\par Esta asserção tambem é absurda. Por quanto, como provámos, a materia não é nem póde ser eterna, mas foi produzida, e por isso não póde ser a causa \emph{primeira} de todas as coisas.\par Note-se a differença entre este e o precedente argumento. Aqui, \textsc{Deus} considera-se como a causa \emph{efficiente} de tudo o que existe; além, \textsc{Deus} considerou-se tambem como causa \emph{final} de todas as coisas, emquanto tudo tende para \textsc{Deus}, como Ente infinitamente intelligivel e amavel.}.

c) Os entes visiveis são todos \emph{contingentes}, isto é, existem mas podiam deixar de existir. Se são \emph{contingentes}, foram produzidos por outro ente; porque o \emph{contingente}, podendo existir e deixar de existir, não contém em si a razão sufficiente da propria existencia. Este outro ente, que produziu o \emph{contingente}, ou é tambem \emph{contingente}, ou é \emph{necessario}. Se é \emph{necessario}, é \textsc{Deus}. Se é \emph{contingente}, foi produzido por outro, e este por outro, até se encontrar inevitavelmente o ente \emph{necessario} ; porque, se não existisse o ente \emph{necessario}, não poderia existir nenhum ente \emph{contingente}, pois repugna uma serie de entes \emph{dependentes} sem um \emph{primeiro}, de que todos dependam. Logo existe \textsc{Deus}\footnote{Em poucas palavras: O mundo é \emph{contingente}. Ora o ente \emph{contingente} depende do ente \emph{necessario}. Logo existe um ente \emph{necessario}. O ente necessario é \textsc{Deus}. Logo existe \textsc{Deus}.\par \emph{Vacherot} affirma muitas vezes (sem nunca o demonstrar) que o mundo,\pgn{435}se é \emph{contingente} na \emph{apparencia}, é \emph{necessario} e \emph{absoluto} na \emph{realidade}. --- Esta affirmação é absurda. O mundo, como demonstrámos na \emph{Cosmologia geral}, é \emph{essencialmente contingente} no seu conjuncto e em cada uma das suas partes.\par \emph{Kant} diz que a prova, deduzida da idea do ser \emph{necessario}, não tem valor; porque, como a prova \emph{ontologica}, passa da ordem \emph{ideal} para a \emph{real}, apoiando-se num principio meramente \emph{subjectivo}, que é o de \emph{causalidade}.\par Rejeitamos com \emph{Kant} a prova \emph{ontologica} ácerca da existencia de \textsc{Deus}, mas negamos que a prova \emph{cosmologica} se reduza á \emph{ontologica}. Na verdade, a prova \emph{cosmologica} não consiste, como pensa \emph{Kant}, em deduzir da \emph{idea} de um ser \emph{necessario} a sua existencia \emph{real}, mas em deduzir da existencia \emph{real} dos seres \emph{contingentes} a existencia \emph{real} do ser \emph{necessario}, que é \textsc{Deus}. --- Além d'isso, negamos que o principio de \emph{causalidade} seja meramente \emph{subjectivo}; elle é tambem, e primeiramente, \emph{objectivo}, porque não é lei do nosso espirito, senão porque é lei das coisas.}.
\pg{435} d) Os entes creados, embora possuam todos uma perfeição \emph{finita}, não a possuem todos no mesmo grau, mas um é \emph{mais} perfeito, e outro o é \emph{menos}. Isto mostra que a perfeição das creaturas é \emph{participada}; porque, se o não fosse, não haveria razão por que numa creatura a perfeição fosse \emph{maior}, e noutra \emph{menor}. Sendo \emph{participada}, a perfeição das creaturas deve derivar de um ente, e de um ente, que a possua pela \emph{sua propria essencia}. Ora a perfeição, que um ente possue pela sua \emph{propria essencia}, deve ser \emph{infinita}; porque, se fosse \emph{finita}, seria \emph{participada}. Logo existe um ente \emph{infinitamente} perfeito, de que derivaram, como de um oceano immenso, as muitas e diversas perfeições das creaturas. O ente \emph{infinitamente} perfeito é \textsc{Deus}. Logo existe \textsc{Deus} \footnote{Em poucas palavras: Os entes creados são todos \emph{finitos}. Ora a existencia dos entes \emph{finitos} suppõe necessariamente a existencia do ente \emph{infinito}. Logo existe o ente \emph{infinito}. O ente \emph{infinito} é \textsc{Deus}. Logo existe \textsc{Deus}.\par Objectam: A diversidade das perfeições, que se encontram nos seres finitos, não demonstra a existencia de um ente infinito.\par A objecção já está refutada no texto. Entre os seres finitos, um é \emph{mais} ou \emph{menos} perfeito que o outro. Ora uma coisa diz-se \emph{mais} ou \emph{menos} perfeita, emquanto \emph{mais} ou \emph{menos} se approxima de um ente \emph{absolutamente perfeito}: assim uma photographia diz-se \emph{mais} ou \emph{menos} perfeita, emquanto \emph{mais} ou \emph{menos} se approxima do original, que é uma coisa \emph{absoluta}. Logo deve existir um ente \emph{absoluta} e \emph{infinitamente} perfeito, que seja a causa das \emph{varias} e \emph{limitadas} perfeições, que se encontram nas creaturas. --- O \emph{relativo}, que se não refira ao \emph{absoluto}, é uma contradicção.}.

e) A intelligencia humana conhece verdades \emph{absolutas} e

\emph{necessarias}, tanto da ordem \emph{especulativa} como da \emph{pratica}. Taes verdades devem ter uma \emph{razão sufficiente}. Esta \emph{razão sufficiente} não se encontra nos \emph{seres creados}, nem na propria \emph{intelligencia humana}. Não se encontra nos \emph{seres creados}; pois estes, sendo \emph{relativos} e \emph{contingentes}, não podem produzir o \emph{absoluto} e o\pg{436} \emph{necessario}. Nem se encontra na \emph{intelligencia creada}; pois esta não forma as verdades \emph{absolutas} e \emph{necessarias}, descobre-as apenas e recebe-as como regra do seu pensamento. Logo a \emph{razão sufficiente} d'aquellas verdades deve encontrar-se numa \emph{intelligencia absoluta} e \emph{necessaria}; pois ellas são entidades \emph{ideaes}, \emph{absolutas} e \emph{necessarias}. A intelligencia \emph{absoluta} e \emph{necessaria} é a de \textsc{Deus}. Logo existe \textsc{Deus} \footnote{E' certo que, além da ordem \emph{physica}, composta de substancias materiaes, existe uma ordem \emph{logica}, composta de verdades \emph{necessarias}, \emph{absolutas}; por quanto estas verdades não podem ser formadas nem mudadas pela nossa intelligencia. Esta faculdade limita-se a descobril-as, e descobre-as como independentes de toda a experiencia e de todo o pensamento creado. Logo ou devemos negar a verdade do principio \emph{de razão sufficiente}, ou deve existir uma intelligencia \emph{necessaria} e \emph{absoluta}, em que estas verdades encontram o seu fundamento. Ora o principio \emph{de razão sufficiente} é verdadeiro. Logo deve existir uma intelligencia \emph{absoluta} e \emph{necessaria}, que é a de \textsc{Deus}.\par Nem se diga que, neste argumento, passamos da ordem \emph{ideal} para a \emph{real}. Por quanto, as verdades \emph{absolutas} e \emph{necessarias} são \emph{objectivas}; pois, como dissemos, ellas são verdades independentemente do nosso pensamento, não podemos formal-as nem mudal-as, são sempre o que são, embora não pensemos nellas e as contestemos. Sendo \emph{objectivas}, devem encontrar na \emph{realidade} um fundamento. Este fundamento não se encontra nas coisas mutaveis, nem na intelligencia creada. Logo deve encontrar-se num Ente \emph{necessario} e \emph{eterno}, em \textsc{Deus}.}.

13. \textbf{Argumentos physicos}. --- Os argumentos \emph{physicos} baseam-se em que, se existe uma ordem maravilhosa no mundo e nas operações dos agentes \emph{irracionaes}, deve existir um ordenador infinitamente sabio e poderoso.

a) Uma ordem maravilhosa e constante existe entre as muitas e diversas substancias, que compõem o universo. Esta ordem exige uma causa infinitamente sabia e poderosa; pois só essa causa, --- por ser \emph{infinitamente sabia}, podia conhecer a natureza de todas as substancias e o posto de cada uma d'ellas no plano da creação, e escolher as leis mais aptas para a constancia e regularidade da mesma ordem; --- e, por ser \emph{infinitamente poderosa}, podia executar esse plano. A causa infinitamente \emph{sabia} e \emph{poderosa} é \textsc{Deus}. Logo existe \textsc{Deus} \footnote{Não insistimos neste ponto; porque na \emph{Cosmologia geral} provámos que a ordem \emph{physica}, com as suas \emph{leis}, é effeito de uma intelligencia \emph{infinita}.\par \emph{Kant} diz: A ordem \emph{physica} é effeito de uma intelligencia, mas não de uma intelligencia \emph{infinita}. Logo não demonstra a existencia de \textsc{Deus}.\par O que dizemos no texto mostra a falsidade da affirmação de \emph{Kant}. --- Além d'isso, a intelligencia, que estabeleceu a ordem no mundo, só podia\pgn{437}ser a do Ente, que \emph{criou} o proprio mundo. Ora o ente, que \emph{criou} o mundo, deve ter um poder \emph{infinito}. O ente, dotado de poder \emph{infinito}, deve ter uma natureza \emph{infinita} e por isso uma intelligencia \emph{infinita}.\par Os materialistas objectam: Esta ordem e regularidade, que existe em todo o universo e em cada substancia, não é effeito de uma causa intelligente e deliberada, mas é obra do \emph{acaso}.\par Resposta. O \emph{acaso}, sendo \emph{privação} de intelligencia e deliberação, é \emph{nada}. Ora é absurdo dizer que o \emph{nada} produziu a \emph{realidade}. Logo a ordem, que existe no mundo, sendo uma \emph{realidade}, não foi produzida pelo \emph{acaso}, mas por um Ente real, infinitamente sabio e poderoso.}.

b) Não só a ordem, que existe no universo, mas tambem a que existe nas operações de cada agente \emph{irracional}, proclama \pg{437} a existencia de uma Intelligencia soberana e infinita. Na verdade, todo o agente \emph{irracional}, nas suas operações, tende para um fim, empregando para isso os meios mais efficazes, e executando um plano com maravilhosa exactidão, constancia e uniformidade. Este fim, esta escolha de meios e este plano, devem ser effeitos de uma intelligencia \emph{infinita}; porque só uma intelligencia \emph{infinita} podia determinar a cada agente \emph{irracional} um fim proporcionado, modelar planos complicados e ao mesmo tempo perfeitos, e, para os realizar, servir-se da forças cegas da natureza. A intelligencia \emph{infinita} é propria de \textsc{Deus}. Logo existe \textsc{Deus} \footnote{Tambem este ponto ácerca da \emph{finalidade} nas operações dos agentes \emph{irracionaes} foi sufficientemente tratado na \emph{Cosmologia geral}.}.

14. \textbf{Argumentos moraes.} --- Os argumentos \emph{moraes} fundam-se em que, --- \emph{a}) se existe uma lei \emph{moral}, deve existir um legislador, --- \emph{b}) se todos os povos admittiram a existencia de um Ente supremo, Este deve existir.

a) A consciencia attesta-nos que temos gravada no nosso coração uma lei moral, á qual não nos podemos subtrahir, e que o cumprimento d'essa lei produz tranquillidade, assim como a sua transgressão causa remorso. Ora a existencia d'esta lei, como a respectiva sancção, suppõe a existencia de um legislador. Este legislador não póde ser o proprio homem; aliás seria inexplicavel o sentimento da nossa dependencia de um poder superior. Logo deve ser um Ente superior ao homem e dotado de virtude ou poder \emph{infinito}; porque só um ente infinitamente poderoso podia gravar no coração dos homens essa lei. O Ente superior, de que o homem depende e a que deve obedecer, é \textsc{Deus}. Logo existe \textsc{Deus} \footnote{Este argumento, apresentado tambem por \emph{Kant} para a demonstração da existencia de \textsc{Deus}, será mais desenvolvido na \emph{Moral}.\par Aqui apenas advertimos que a existencia da \emph{lei moral} considera-se como um \emph{facto experimental}, attestado pela consciencia. Provado o facto da existencia da \emph{lei}, legitimamente deduz-se a existencia de um \emph{legislador}.}.
\pg{438} b) Todos os homens, em todos os tempos e em todos os logares, reconheceram e admittiram a existencia de um Ente supremo. Este consenso unanime, constante e universal, não póde ser effeito do \emph{erro}; pois o \emph{erro} é causa de discordias, é mutavel e particular. Logo é effeito da \emph{verdade}; porque só a \emph{verdade} é unanime, constante e universal, e só ella pode levar os homens a mortificar as paixões e renunciar á propria vontade, para se sujeitarem á vontade de um Ente supremo. Logo existe \textsc{Deus} \footnote{O consenso do genero humano, relativamente á \emph{existencia} de um Ente supremo, não podia ser mais constante, nem mais universal. Esta crença num Ente supremo podia ser e foi desfigurada por muitos erros; mas conservou sempre o fundo intacto. Muitos povos ignoraram e ignoram ainda a \emph{essencia} de \textsc{Deus}, mas nenhum ignorou ou ignora a sua \emph{existencia}.\par \emph{Cicero} disse: « Não existe povo algum, por inculto e selvagem, que não tenha fé em \textsc{Deus}, embora lhe não conheça a essencia » (\emph{De nat}. \emph{Deor}. I, 17).\par \emph{Quatrefages} escreve: « Obrigado pelo meu ensino a passar revista a todas as raças humanas, procurei o \emph{atheismo} tanto entre os povos mais selvagens, como entre os mais civilisados. \emph{Não o encontrei em parte alguma, a não ser nalgum individuo ou eschola muito limitada}, como se viu na Europa no seculo passado e como se vê ainda no presente » (\emph{L'espèce humaine}, pag. 355, Paris, 1877)... E' noutro logar não se exprime menos claramente: « Pouco a pouco fez-se a luz, e os povos da \emph{Australia} e da \emph{Melanesia}, os \emph{Boschimans}, os \emph{Ottentotes}, os \emph{Cafres}, os \emph{Bechuanas}, foram eliminados do numero dos povos atheus e reconhecidos como religiosos » (\emph{Rapport}, pag. 410).\par De resto, os \emph{Annaes da propagação da Fé} provam, com razões claras e baseadas na experiencia, que povos atheus não passam de um mytho.\par Objectam: A crença de todos os povos na existencia de um Ente supremo derivou da \emph{ignorancia}, ou da astucia dos \emph{sacerdotes}, ou da prepotencia dos \emph{legisladores}, ou do \emph{medo}, ou dos \emph{prejuizos}.\par Resposta: A crença unanime de todos os povos em \textsc{Deus} não derivou --- \emph{a}) nem da \emph{ignorancia}, --- \emph{b}) nem da astucia dos \emph{sacerdotes}, --- \emph{c}) nem da prepotencia dos \emph{legisladores}, --- \emph{d}) nem do \emph{medo}, --- \emph{e}) nem dos \emph{prejuizos}.\par a) Não derivou da \emph{ignorancia}. Por quanto, o estudo das causas naturaes é o meio mais facil e efficaz para se conhecer a existencia de \textsc{Deus}; visto que essas causas, sendo \emph{secundarias}, \emph{contingentes}, \emph{imperfeitas}, exigem uma causa \emph{primeira}, \emph{necessaria}, \emph{perfeitissima}, que é \textsc{Deus}.\par A experiencia confirma a these. Na verdade, os maiores naturalistas são os mais convencidos e enthusiastas adoradores de \textsc{Deus}. --- Citaremos apenas os testemunhos de \emph{Linneu} e \emph{Kepler}.\par \emph{Linneu} escreve: « Quando acordei, \textsc{Deus} \emph{eterno}, \emph{immenso}, \emph{omnisciente}, \emph{omnipotente}, acabava de passar. Vi-o ao longe, e fiquei profundamente maravilhado. Segui os vestigios de seus passos, atravez das obras da creação, e por toda a parte, mesmo nas mais pequenas coisas, que poder! que sabedoria! que ineffavel perfeição! O sol e todo o systema planetario, immenso, incalculavel, appareceram-me suspensos pelo Primeiro Motor, Causa das\pgn{439}causas, Governador e Conservador do universo. Todas as coisas são, pois, testemunho da sabedoria divina; a belleza, a harmonia, as justas proporções das creaturas proclamam o poder do grande \textsc{Deus} » (\emph{Syst}. \emph{de la Nat}. p. 10).\par \emph{Kepler} conclue a sua obra, intitulada « \emph{Os quatro livros da harmonia celeste} », por estas sublimes palavras: « E agora só me resta levantar as mãos e os olhos para ao ceu, e dirigir ao Auctor de toda a luz a minha humilde oração: O' Vós, que, pelos esplendores sublimes que derramastes sobre toda a natureza, elevais os nossos desejos até á divina luz da vossa graça, sede bemdito! Senhor e Creador de todas as consolações, que experimentei no extasis, a que me elevou a contemplação da obra das vossas mãos, eis acabado este livro, que contém o fructo do meu trabalho e de toda a intelligencia, que me destes. Proclamei deante dos homens toda a grandeza das vossas obras, mostrando-lhes, na medida das minhas forças, a sua immensa perfeição. Esforcei-me por chegar até á verdade, por a conhecer o mais perfeitamente possivel, e, se me escapou alguma coisa indigna de Vós, fazei que a conheça para a eliminar. Ter-me-hei deixado seduzir pela presumpção, na presença da belleza maravilhosa das vossas obras? Terei eu procurado a minha gloria entre os homens, levantando este monumento, que só devia ser consagrado á vossa honra? Ah! se assim foi, tende piedade e misericordia de mim, e concedei-me a graça de que o meu trabalho não produza mal algum, mas concorra para a glorificação do vosso Nome e para a salvação das almas ».\par Os mesmos sentimentos encontram-se nos escriptos dos mais celebres naturalistas, como \emph{Copernico}, \emph{Galileu}, \emph{Newton}, \emph{Jussieu}, \emph{Cuvier}, \emph{Agassiz}, \emph{Biot}, \emph{Buffon}, \emph{Lavoisier}. \emph{Brongniart}, \emph{Blainville}, \emph{Bertholet}, \emph{Gay-Lussac}, \emph{Elie de Beaumont}, \emph{Thénard}, \emph{Réaumur}, \emph{Barande}, \emph{Milne-Edwards}, \emph{Dumas}, \emph{Chevreul}. \emph{Liebig}, \emph{Flourens}, \emph{Quatrefages}, \emph{Pasteur}, etc. Nem admira; porque, como diz \emph{Herschel}, « quanto mais se extende o campo da Physica, tanto mais numerosas e\pgn{440}inconcussas se tornam as demonstrações da existencia eterna de uma intelligencia creadora e omnipotente; os geologos, os astronomos, os naturalistas trouxeram todos a sua pedra para este grande templo da sciencia, --- templo levantado ao proprio \textsc{Deus} ». --- Não nos cançaremos de repetir com \emph{Bacon}; « A pouca sciencia afasta de \textsc{Deus}, a muita conduz a Elle ».\par b) Não derivou da astucia dos \emph{sacerdotes}; pois a instituição dos \emph{sacerdotes é effeito}, e não \emph{causa}, da crença universal na Divindade.\par c) Não derivou da prepotencia dos \emph{legisladores}. Os legisladores basearam as leis na propria crença religiosa dos seus povos. --- Além d'isso, os povos não teriam acceitado, ao menos em toda a parte, uma crença, que é contraria a todas as paixões. --- Finalmente, as leis variam em todos os povos, mas a crença em \textsc{Deus} foi sempre a mesma por toda a parte.\par d) Não derivou do \emph{medo}. Na verdade, esse \emph{medo} ou era conforme com a razão, ou não. Se não era conforme com a razão, não podia produzir a crença em \textsc{Deus}; pois essa crença é universal e constante, e um effeito universal e constante não póde derivar de uma causa particular e mutavel, como é o \emph{medo}. Se era conforme com a razão, esse \emph{medo} já suppunha a crença em \textsc{Deus}; pois ninguem tem \emph{medo} de um ente, que se não conhece como existente. --- A verdade é que o \emph{medo}, longe de produzir a crença em \textsc{Deus}, deu sempre occasião ao \emph{atheismo}, pois os libertinos, \emph{pelo facto de temer o castigo}, esforçam-se por negar a existencia de \textsc{Deus}.\par e) Não derivou dos \emph{prejuizos}; pois os \emph{prejuizos} variam segundo a diversidade dos povos, ao passo que a crença em \textsc{Deus} é universal e constante.}.
\pg{439}

%%>>> art-38

\section*{Artigo III.\\ A existencia de Deus e o atheismo.}
\addcontentsline{toc}{section}{Artigo III. --- A existencia de Deus e o atheismo.}

15. \textbf{Atheismo.} --- \emph{Atheismo} é o erro dos que não reconhecem ou negam a \emph{existencia} de \textsc{Deus}.

16. \textbf{Divisão do atheismo.} --- O \emph{atheismo} divide-se em \emph{pratico} e \emph{theoretico}. --- O \emph{pratico} é o erro dos que, embora affirmem a existencia de \textsc{Deus} com a \emph{razão}, negam-na com os \emph{factos}, pois vivem, como se \textsc{Deus} não existisse. --- O \emph{theoretico} é o erro dos que negam com a \emph{razão} a existencia de \textsc{Deus}.

17. \textbf{Divisão do atheismo theoretico.} --- O atheismo \emph{theoretico} póde ser \emph{negativo} ou \emph{positivo}. --- O \emph{negativo} é o erro dos que não admittem a existencia de \textsc{Deus}, porque o não conhecem. --- O \emph{positivo} é o erro dos que, tendo conhecido a \textsc{Deus}, estão \emph{séria} e \emph{sinceramente} convencidos de que Elle não existe.
\pg{440} 18. \textbf{O atheismo pratico é possivel.} --- O atheismo \emph{pratico} não é uma depravação da \emph{intelligencia,} mas sim da \emph{vontade.} Ora a nossa vontade está sujeita á depravação, e, para se entregar mais livremente ás desordens, póde deixar de pensar em \textsc{Deus} e de Lhe prestar as homenagens da sua piedade e obediencia. --- Este \emph{atheismo} não só é \emph{possivel,} mas existe realmente\footnote{\emph{Bayle} diz: « Os que não servem a \textsc{Deus}, e que antes o offendem com uma vida cheia de crimes e vicios, são uma especie de atheus, é o atheismo dos libertinos e devassos; porque effectivamente vivem, como se não houvesse \textsc{Deus}, sem medo da sua Justiça, sem gratidão para com a sua Bondade, sem respeito para o seu Nome, sem obediencia ás suas leis, e, ainda que podessem tirar \textsc{Deus} do seu throno para o esmagar sob os seus pés e aniquilar para sempre, não viveriam com mais licença e ousadia. São atheus pelas suas acções; conhecem a \textsc{Deus} com a lingua, mas negam-no com os costumes; fallam como crentes, mas operam como incredulos » (\emph{Pensées div.,} t. IV, pag. 94).}.

19. \textbf{O atheismo theoretico-negativo não é possivel nos homens adultos.} --- A existencia de \textsc{Deus}, embora nos não seja \emph{immediatamente} evidente, todavia descobre-se logo que começamos a applicar os primeiros principios da razão. Por quanto, sendo natural e facillimo deduzir da existencia dos effeitos a existencia da causa respectiva, é claro que, contemplando o universo, não podemos deixar de reconhecer uma causa, que o\pg{441} produziu e formou. Logo o atheismo \emph{theoretico-negativo} é impossivel nos homens \emph{adultos} \footnote{A historia contém exemplos de homens que, isolados, desde crianças, de toda a convivencia social, chegaram, só pela força da razão, ao conhecimento de uma causa productora do mundo.}.

20. \textbf{O atheismo theoretico-positivo é impossivel.} --- E' impossivel que o homem esteja \emph{séria} e \emph{sinceramente} convencido da falsidade de uma coisa, attestada pela voz da natureza. Ora a existencia de \textsc{Deus} é attestada pela voz da natureza; pois é admittida por todos os homens, em todos os tempos, em todos os logares, e só o que deriva da natureza se encontra em todos os homens, em todos os tempos, em todos os logares. Logo é impossivel que o homem esteja \emph{séria} e \emph{sinceramente} convencido de que não existe \textsc{Deus}. --- A experiencia confirma a these; pois os que se dizem \emph{atheus} são \emph{scepticos} ou \emph{pantheistas} \footnote{O nosso seculo foi chamado \emph{era do atheismo}. E' certo que alguns infelizes empregam esforços para destruir, se fosse possivel, a propria idea de \textsc{Deus}. \emph{Fouerbach} diz: « Só o homem é o verdadeiro salvador! Só o homem é o nosso \textsc{Deus}, o nosso Juiz, o nosso Redemptor ». \emph{Rénan} escreve: « \textsc{Deus}, a Providencia, a immortalidade, não passam de palavras velhas e pesadas ». A loja maçonica de Liege, em 1865, dizia: « Não só devemos collocar-nos acima de todas as religiões, mas acima de toda a crença num Deus ».\par Mas as pretendidas \emph{convicções} d'esses desvairados não passam de \emph{duvidas}. O espirito nunca poderá sinceramente convencer-se de que \textsc{Deus} não existe e, pelo menos com as suas duvidas, protestará sempre contra as estultas affirmações do coração. Já nos seus tempos \emph{Seneca} dizia: « Mentem os que dizem estar convencidos de que \textsc{Deus} não existe; pois, se ás vezes o dizem, todavia de noite, quando estão sós, começam a duvidar ».\par Estas \emph{duvidas} ácerca da existencia de \textsc{Deus} derivam da \emph{superficialidade} do pensamento e, sobretudo, da \emph{corrupção} da vontade.\par As \emph{duvidas} derivam da \emph{superficialidade} do pensamento. A este respeito, \emph{Caro} escreve: « Ha entre nós uma classe de espiritos scientificamente desorganisados, habeis para comprehender os phenomenos e as leis, nas diversas ordens da realidade, na natureza ou na historia, mas incapazes, por causa de uma inercia systematica, que se toma por uma força, de se elevar á concepção da \emph{causa}, que é a razão ultima das leis naturaes, assim como as leis são a razão dos factos » (\emph{L'idée de Dieu}, 1873, pag. 42).\par Mas as \emph{duvidas} derivam principalmente da \emph{corrupção} da vontade. Não é na \emph{intelligencia}, que o atheu nega a existencia de \textsc{Deus} ou duvida d'ella, mas é no \emph{coração}. O homem, cuja vida é pura e ordenada, nem sequer duvida da existencia de \textsc{Deus}. \emph{Rousseau} dizia: « Conservai a vossa alma num estado de desejar que exista \textsc{Deus}, e nunca duvidareis da sua existencia ». \emph{La Bruyère} escreve: « Desejava ouvir a um homem sabio, moderado, casto e justo, que \textsc{Deus} não existe ou que a nossa alma não é immortal; elle fallaria sem interesse. Mas esse homem não se encontra ».}

21. \textbf{Consequencias do atheismo.} --- O \emph{atheismo} é a destruição da \emph{sciencia}, da \emph{sociedade} e da \emph{natureza humana}.
\pg{442} a) E' a destruição da \emph{sciencia}. Na verdade, a sciencia é o conhecimento das \emph{causas} dos seres. Ora, se \textsc{Deus} não existisse, os problemas mais fundamentaes seriam inexplicaveis. Porquanto, --- na ordem \emph{physica,} seria inexplicavel a origem do mundo, da vida, da alma humana, a uniformidade e constancia da ordem e das leis physicas; --- na ordem \emph{logica,} faltaria a base da \emph{possibilidade} das coisas, da \emph{eternidade} e \emph{necessidade} das essencias, da \emph{immutabilidade} da verdade; --- na ordem \emph{moral,} seria inexplicavel a existencia da \emph{lei moral,} a virtude não teria o premio, nem o vicio o castigo, e o bem não differiria do mal.

b) E' a destruição da \emph{sociedade}. Porquanto, uma sociedade não póde subsistir, sem que haja fidelidade e respeito nos subditos, desinteresse e justiça no principe, amor e concordia em todos. Mas uma sociedade sem \textsc{Deus} é o reino da anarchia, tendo por direito a força material, e por lei o despotismo (1).\footnote{\emph{Voltaire} diz (\emph{Dict. phil.}, a. \emph{athée}): « Eu não quereria ser subdito de um principe atheu, que, se encontrasse o seu interesse em me esmigalhar, com certeza esmigalhar-me-hia. Eu não quereria, se fosse soberano, ter cortezãos atheus, pois, se tivessem interesse em me envenenar, deveria tomar, todos os dias, contra-venenos. E', pois, absolutamente necessario, para os principes e para os povos, que a idea de um Ente supremo, creador, governador, remunerador, seja gravada profundamente em todos os espiritos ».}

c) E' a destruição da \emph{natureza humana}. Com effeito, se \textsc{Deus} não existisse, a nossa alma não seria immortal, e assim o homem seria a mais desgraçada das creaturas; porque, depois de ter soffrido tantas dores, depois de ter suspirado tanto pela felicidade sem nunca a encontrar, o \emph{nada} seria a recompensa dos seus soffrimentos e o descanço das suas tendencias (2).\footnote{Estas consequencias do \emph{atheismo} são tambem provas, ainda que \emph{indirectas}, da nossa these. Comtudo, \emph{Stuart Mill} diz que não temos prova alguma, nem \emph{seriamente provavel}, para demonstrar a existencia de \textsc{Deus}!}
